\subsection{What critical amplification would force}
\label{sec:critical-balance}

The coefficient $1/5$ comes from balancing three upper bounds.
Near equality therefore forces several structural conditions
simultaneously. The following statements concern a normalized circuit
with $q\ge1$ and a witness-dependent gate as output, using the notation
of Section~\ref{sec:consistency-graph}.

\begin{lemma}[Finite stability of the three-way balance]
\label{lem:finite-stability}
Fix $\delta>0$, put $\alpha=1/3+\delta$, and choose $A_\delta\ge1$
large enough for
\[
 \C(g)\le p+A_\delta(q+1)2^{\min\{h,b,\alpha\kappa\}}.
\]
If $t\ge0$ and $\C(g)\ge p+A_\delta(q+1)2^t$, then
\[
 \Delta=q+1-(2+1/\alpha)t\ge0,
\]
and there are $e_h,e_b,e_\kappa\ge0$ satisfying
\[
 h=t+e_h,\qquad b=t+e_b,\qquad
 \kappa=t/\alpha+e_\kappa,\qquad
 e_h+e_b+e_\kappa\le\Delta.
\]
\end{lemma}
\begin{proof}
Comparing the two size inequalities forces
$h,b,\alpha\kappa\ge t$. Subtract these lower bounds from
$h+b+\kappa\le q+1$.
\end{proof}

This elementary lemma is useful because it preserves the fixed-$\delta$
constant before any limit is taken. Its asymptotic consequence also
constrains the residual graph, beyond the original cycle count.

\begin{theorem}[Necessary structure at the critical gate budget]
\label{thm:critical-structure}
Consider a family with $m\to\infty$ declared witness bits, normalized
size $s=p+q$, and
\[
 s+1=(5+o(1))m,\qquad \log\C(g)\ge m-o(m).
\]
Let $a$ be the number of used ordinary inputs. Then
\begin{align*}
 p&=o(m),& q&=(5+o(1))m,\\
 a,h,b&=(1+o(1))m,& \kappa&=(3+o(1))m.
\end{align*}
For every choice of legal equality splitting and reductions in
Lemma~\ref{lem:cubic-reduction}, the resulting kernel $K$ is
eventually nonempty and satisfies
\begin{align*}
 |V(K)|&=(6+o(1))m,& \kappa-\kappa(K)&=o(m),\\
 \operatorname{cw}(K)&=(1+o(1))m,&
 \operatorname{pw}(K)&=(1+o(1))m.
\end{align*}
Here cutwidth and pathwidth mean their minimum over all orders and
decompositions. No hard family satisfying the hypotheses is asserted.
\end{theorem}
\begin{proof}
Since $p=O(m)$ whereas $\C(g)\ge2^{m-o(m)}$,
$\log(\C(g)-p)\ge m-o(m)$. For each fixed $\delta>0$,
Theorem~\ref{thm:fifth} therefore gives
\[
 h\ge m-o(m),\qquad b\ge m-o(m),\qquad
 \kappa\ge m/\alpha-o(m).
\]
The prefactor contributes only $O_\delta(\log m)$ to the logarithm.
Using $p+h+b+\kappa\le s+1$, we obtain
\[
 \limsup \frac{p}{m}\le3-1/\alpha,\qquad
 \limsup \frac{h}{m},\ \limsup\frac{b}{m}\le4-1/\alpha.
\]
These inequalities hold for every fixed $\delta$. Letting $\delta$
tend to zero after taking the limits proves $p=o(m)$ and
$h/m,b/m\to1$. The budget and the fixed-$\delta$ lower bounds
then give $\kappa/m\to3$, and $q/m\to5$ follows from $s=p+q$.

To relate $a$ to $h$, note first that $h\le a+p$. In the undirected
graph on the $a$ ordinary inputs and $p$ independent gates, every
connected component contains an interface signal, because every
vertex is in the output cone. If there are $c$ components, then
$c\le h$, and counting edges gives
$a+p-c\le2p$. Thus $a\le p+h$, so $|a-h|\le p$ and $a/m\to1$.

Fix any sequence of legal reductions along the family. An empty
kernel would give $\C(g)=O(s)$, so eventually it is nonempty.
Write $N=|V(K)|$. The finite compiler bound, applied to an order
of minimum cutwidth, gives
\[
 \C(g)\le p+400(q+1)2^{\operatorname{cw}(K)},
 \qquad \operatorname{cw}(K)\ge m-o(m).
\]
The cutwidth lower bound and $\operatorname{cw}(K)\le3N/2$ imply
$N\to\infty$. On the other hand $N\le2(\kappa-1)\le6m+o(m)$. For every fixed
$\eta>0$, the cubic pathwidth theorem and
Lemma~\ref{lem:median-width} give, for sufficiently large $N$,
\[
 \operatorname{cw}(K)\le\operatorname{pw}(K)+2
       \le(1/6+\eta)N+2.
\]
Hence $\liminf N/m\ge1/(1/6+\eta)$ for every $\eta>0$, proving
$N/m\to6$. The cubic identity $\kappa(K)=N/2+1$ gives the cycle
loss. The same inequalities sandwich each width between $m-o(m)$
and $m+o(m)$. All constants depend only on the fixed slack, so this
argument applies to any selected legal reduction sequence.
\end{proof}

There is little unused fanin capacity in such a family. If $r$ counts
independent-signal input pins of the dependent gates, with multiplicity,
and $U=2q-(\ell+r)$, then
\[
 q+1-h-b-\kappa=U+(r-h)=o(m).
\]
Both terms are nonnegative. Thus only $o(m)$ input slots are missing
from binary gates, and ordinary signals have only $o(m)$ repeated
uses at the interface. This concerns distinct signals, not semantic
independence of their values.

\paragraph{Combining the structural and witness conditions.}
The hypotheses hold whenever $s+1=(5+o(1))m$ and
$\C(g)\ge T=2^m(s+1)/L$ with $\log L=o(m)$.
Corollary~\ref{cor:sparse-population} then also forces
$\Omega((T/n)\log(2+T/n))$ positive ordinary inputs with $O(nL)$
declared witnesses. Removing the $m-b$ unused witness bits divides
each positive fiber and this threshold by exactly $2^{m-b}$;
the two witness conventions must be kept consistent.
Enumeration already gives $m-b\le\log L$ under this amplification
hypothesis. These are simultaneous necessary conditions. Large
width, balanced budgets, and sparse fibers alone do not establish
a lower bound on $\C(g)$.
